\textbf{COLUNI 2021 -- Questão 08}\par\medskip De acordo com a variante padrão da Língua Portuguesa, o poeta, romancista e dramaturgo brasileiro Oswald de Andrade cometeu um deslize gramatical ao escrever o verso “Me dá um cigarro”, em seu poema \textit{Pronominais}.

\par

Assinale a alternativa em que a narradora Fortunatella, do livro \textit{A menina que fez a América}, de Ilka B. Laurito, comete o mesmo deslize que o poeta:\begin{enumerate}
\item “Lá me parecia um lugar assombrado e perigoso.” (LAURITO, Ilka Brunhilde. \textbf{A Menina que fez a América}. São Paulo,
      FTD, 2019. p. 37)
\item “Ela não usou a palmatória, mas é como se tivesse me batido duramente.” (LAURITO, Ilka Brunhilde. \textbf{A Menina que fez a América}. São Paulo, FTD, 2019. p. 60)
\item “Um dia meus primos me convidaram para passear de burrico com eles e a prima Rina.” (LAURITO, Ilka
      Brunhilde. \textbf{A Menina que fez a América}. São Paulo, FTD, 2019. p. 43)
\item “Ela talvez me amasse, pois não enviava apenas frias libras esterlinas, mas alguma coisa bonita e
      especial.” (LAURITO, Ilka Brunhilde. \textbf{A Menina que fez a América}. São Paulo, FTD, 2019. p. 68)\end{enumerate}
